\documentclass[11pt]{article}
\usepackage{amsmath,amssymb,amsthm,hyperref}
\newtheorem{theorem}{Theorem}
\newtheorem{proposition}[theorem]{Proposition}
\newtheorem{corollary}[theorem]{Corollary}
\title{Boolean Projection Identities and Algebraic Switching}
\author{}
\date{November 23, 2026}
\begin{document}
\maketitle

\paragraph{Scope.}
This note gives an exact algebraic description of simultaneous
integer-affine multiset cancellations on the Boolean cube. It also
identifies the compatibility conditions for prescribed polynomial
projections. These are manual results. They do not assert that the
required cancellations have small support, or that arbitrary Boolean
signs admit the requested head count.

\section{The projection kernel}

Let $R$ be an integral domain, let $n\geq 0$, and set
\[
 \mathcal A=R[t,z_1,\ldots,z_n],\qquad
 J=((z_i-1)(z_i-t):1\leq i\leq n).
\]
For $T\subseteq[n]$, define the $R[t]$-algebra map
\[
 \phi_T:\mathcal A\longrightarrow R[t],\qquad
 \phi_T(z_i)=\begin{cases}t&i\in T,\\1&i\notin T.\end{cases}
\]
Write $z_S=\prod_{i\in S}z_i$ and
$w_S=\prod_{i\in S}(z_i-1)$, with empty products equal to one.

\begin{theorem}[Kernel and canonical remainder]
Every $F\in\mathcal A$ has a unique remainder modulo $J$ of the form
\[
 \overline F=\sum_{S\subseteq[n]}a_S(t)z_S
            =\sum_{S\subseteq[n]}c_S(t)w_S.
\]
Moreover,
\[
 J=\bigcap_{T\subseteq[n]}\ker\phi_T.
\]
In particular $J$ is radical, and $\mathcal A/J$ is a free
$R[t]$-module of rank $2^n$.
\end{theorem}
\begin{proof}
Divide successively in each variable by the monic polynomial
$z_i^2-(1+t)z_i+t$. The resulting remainder has degree at most one
in each $z_i$; each division changes the polynomial by an element
of $J$. The change from the basis $z_S$ to the basis $w_S$ is
triangular with diagonal entries one, so either displayed form
exists over $R[t]$.

For a remainder in the second form, evaluation gives
\[
 \phi_T(\overline F)=\sum_{S\subseteq T}(t-1)^{|S|}c_S(t).
\]
M\"obius inversion on subsets therefore gives
\[
 (t-1)^{|S|}c_S(t)
 =\sum_{T\subseteq S}(-1)^{|S|-|T|}\phi_T(\overline F).
\]
If all evaluations vanish, the right side is zero. Since $R[t]$
is an integral domain and $t-1\neq0$, every $c_S$ is zero.
This proves both uniqueness and the asserted kernel equality:
each generator of $J$ vanishes under every $\phi_T$, and a
polynomial with all evaluations zero has zero remainder.
The remainder monomials consequently give the asserted free basis.
Each $\ker\phi_T$ is prime because its quotient is the integral
domain $R[t]$. Their intersection is radical.
\end{proof}

\begin{theorem}[Exact compatibility conditions]
Given polynomials $G_T(t)\in R[t]$ for all $T\subseteq[n]$, there
exists $F\in\mathcal A$ such that $\phi_T(F)=G_T$ for every $T$
if and only if, for every $S\subseteq[n]$,
\[
 (t-1)^{|S|}\ \bigm|\
 D_S(t):=\sum_{T\subseteq S}(-1)^{|S|-|T|}G_T(t).
\]
When these conditions hold, its unique remainder is
\[
 \overline F=\sum_{S\subseteq[n]}
       \frac{D_S(t)}{(t-1)^{|S|}}w_S.
\]
\end{theorem}
\begin{proof}
The preceding inversion formula proves necessity. Conversely the
divisibility assumptions make every displayed coefficient a polynomial.
For a subset $U$, its evaluation is $\sum_{S\subseteq U}D_S$.
Substitute the definition of $D_S$ and interchange the finite sums.
The coefficient of $G_T$ is
$\sum_{T\subseteq S\subseteq U}(-1)^{|S|-|T|}$, which is one
if $T=U$ and zero otherwise. Thus the evaluation equals $G_U$.
\end{proof}

\paragraph{Compatibility is not automatic.}
Already for $n=1$, the prescriptions $G_\varnothing=0$ and
$G_{\{1\}}=1$ are impossible: their difference is not divisible
by $t-1$. At $t=1$ all projection maps coincide. More generally
the theorem specifies higher-order constraints at that collision.
After adjoining $(t-1)^{-1}$, the displayed formulas interpolate
arbitrary prescriptions, but their coefficients need not be
polynomials before that localization.

\section{Integer-affine lists}

For an integer vector $a=(a_0,\ldots,a_n)$ let
\[
 L_a(x)=a_0+\sum_{i=1}^n a_ix_i,\qquad
 Z^a=t^{a_0}\prod_{i=1}^n z_i^{a_i}.
\]
Negative exponents are permitted in the Laurent polynomial ring
$R[t^{\pm1},z_1^{\pm1},\ldots,z_n^{\pm1}]$.
Every $\phi_T$ extends to this ring, with values in $R[t^{\pm1}]$.

\begin{corollary}[Multiset cancellation]
For two finite lists of integer-affine forms $L_{a_h}$ and
$L_{b_\ell}$, the multisets
\[
 \{L_{a_h}(1_T):h\}\quad\hbox{and}\quad
 \{L_{b_\ell}(1_T):\ell\}
\]
agree for every $T\subseteq[n]$ if and only if
\[
 \sum_h Z^{a_h}-\sum_\ell Z^{b_\ell}
\]
belongs to the Laurent extension of $J$, where the coefficients
are taken in $\mathbb Z$.
\end{corollary}
\begin{proof}
Each projected Laurent polynomial is the difference of the generating
polynomials for the two integer multisets. Its coefficient at an
integer power is the difference of the corresponding multiplicities,
so it vanishes exactly when the multisets agree.
For the kernel assertion, multiply the Laurent polynomial by a
monomial with sufficiently large nonnegative exponents to obtain an
ordinary polynomial. Every projected multiplier is a nonzero power
of $t$. The ordinary polynomial has zero projections exactly when
the Laurent polynomial does. Apply the kernel theorem and then
invert the multiplier.
\end{proof}

The generator $(z_i-1)(z_i-t)$ describes the four-term identity
\[
 \{2x_i,1\}=\{x_i,1+x_i\}
 \qquad(x_i\in\{0,1\}).
\]
Multiplying by a Laurent monomial adds the same affine form to
all four entries. An ideal expression is a finite signed combination
of these translations. It need not preserve distinct entries or
positivity of intermediate multiplicities, and expanding its
coefficient polynomials can increase support. Therefore this algebraic
generation statement is not a bounded-head compilation theorem.

\section{A dense construction and its cost}

Let $r:2^{[n]}\to R$ be arbitrary and define its subset coefficients
\[
 \widehat r(S)=\sum_{T\subseteq S}(-1)^{|S|-|T|}r(T).
\]
For an integer $k\geq n$, set
\[
 F_r(t,z)=\sum_{S\subseteq[n]}\widehat r(S)
       (t-1)^{k-|S|}w_S.
\]
The compatibility theorem, or direct subset inversion, proves
\[
 \phi_T(F_r)=(t-1)^k r(T).
\]
This is an explicit construction for every table. It contains up to
$2^n$ subset terms before expansion. Each factor $(t-1)^{k-|S|}$
has at most $k-|S|+1$ monomials and $w_S$ has $2^{|S|}$ monomials.
Consequently the expansion has at most
\[
 \sum_{S\subseteq[n]}(k-|S|+1)2^{|S|}
  =(k+1)3^n-2n3^{n-1}
\]
monomial occurrences for $n\geq1$. Cancellations can reduce this
number. After collecting terms, the total degree is at most $k$
and each $z_i$ has degree at most one. The number of possible
distinct monomials is consequently at most
\[
 \sum_{S\subseteq[n]}(k-|S|+1)
       =(k+1)2^n-n2^{n-1}.
\]
Even this improved count does not have order $2^n/n^2$. The
unexpanded factorization cannot be charged as one attention head.

For $k=n$, an equivalent expression is
\[
 F_r(t,z)=\sum_{T\subseteq[n]}r(T)
        \prod_{i\in T}(z_i-1)\prod_{i\notin T}(t-z_i).
\]
Only the summand with $T=U$ survives under $\phi_U$, and its
value is $(t-1)^nr(U)$. Both expressions are multilinear in
the $z_i$, so the uniqueness theorem proves their equality.

\begin{proposition}[A positive-integral head compiler]
Work over $\mathbb R$. Suppose a finite Laurent polynomial
\[
 F(t,z)=\sum_{h=1}^H c_h t^{a_{h0}}\prod_i z_i^{a_{hi}}
\]
has integer exponents and satisfies
$y(x)F(t,t^{x_1},\ldots,t^{x_n})>0$ for every Boolean $x$ and
every $0<t<1$, where $y(x)\in\{-1,1\}$.
There is an admissible score with $H$ heads and strict signs $y$.
\end{proposition}
\begin{proof}
Choose an affine function $C(x)=c_0+\sum_i c_ix_i$ such that
every coefficient of every
$B_h(x)=C(x)+a_{h0}+\sum_i a_{hi}x_i$ is strictly positive.
This is possible because the list is finite: for each coordinate,
choose the corresponding coefficient of $C$ larger than one plus
the largest absolute exponent in that coordinate.
Termwise integration of the finite sum gives
\[
 \sum_{h=1}^H\frac{c_h}{B_h(x)}
   =\int_0^1 t^{C(x)-1}F(t,t^{x_1},\ldots,t^{x_n})\,dt.
\]
Each term is integrable because $B_h(x)>0$. The weight
$t^{C(x)-1}$ is positive. The assumed strict sign of the integrand
therefore gives the required strict sign of the integral.
All numerator constants and all denominator coefficients satisfy
the manuscript normal form.
\end{proof}

Applying the proposition to $(-1)^nF_r$ for a nonzero table gives
its signs, with at most $(n+2)2^{n-1}$ terms when $n\geq1$.
The integral can also be evaluated exactly:
\[
 \int_0^1t^{C(x)-1}(1-t)^n\,dt
   =\frac{n!}{C(x)(C(x)+1)\cdots(C(x)+n)}.
\]
For completeness this follows by integration by parts: the value
for $n=0$ is $1/C(x)$, and for $n\geq1$ the integral equals
$n/C(x)$ times the corresponding integral with exponent $C(x)$
and power $n-1$. Iterating gives the formula. This construction
is weaker than the manuscript's already established universal
bound and is recorded only to charge the algebraic support cost
correctly. No compression theorem is hidden in the integral.

\section{The first variation of an ideal product identity}

Let $m\geq1$. Suppose that the affine lists
$u_1(x),\ldots,u_m(x)$ and $v_1(x),\ldots,v_m(x)$ both permute
$0,\ldots,m-1$ at every Boolean input. Let $a_h,b_\ell$ be
arbitrary affine perturbations, and define
\[
 F_\epsilon(s,x)=\prod_{h=1}^m(s+u_h(x)+\epsilon a_h(x)),\qquad
 G_\epsilon(s,x)=\prod_{\ell=1}^m(s+v_\ell(x)+\epsilon b_\ell(x)).
\]
Write $d_j=(-1)^j j!(m-1-j)!$. Then the condition
\[
 \left.\frac{\partial}{\partial\epsilon}
        (F_\epsilon-G_\epsilon)\right|_{\epsilon=0}=r(x),
 \quad\hbox{independently of }s,
\]
is equivalent, separately at every cube input, to
\[
 \sum_h a_h(x)t^{u_h(x)}-\sum_\ell b_\ell(x)t^{v_\ell(x)}
       =\frac{r(x)}{(m-1)!}(1-t)^{m-1}.
\]
Indeed, evaluating the product variation at $s=-j$ gives
$d_j(a_{u^{-1}(j)}-b_{v^{-1}(j)})=r$. Conversely these $m$
values determine the degree-at-most-$m-1$ variation polynomial.
Finally $1/d_j=(-1)^j\binom{m-1}{j}/(m-1)!$, so the binomial
theorem gives the generating-polynomial equation.

The factors $a_h(x)$ and $b_\ell(x)$ are functions of the input.
They are not coefficients independent of $x$ in the preceding
projection ring. To encode them there, one may replace $x_i$ by
$(z_i-1)/(t-1)$, at the price of a localization, and then clear
the factor $t-1$. This observation does not establish sparse
representations with the required affine supports.

\section{Literature check and a Chinese-remainder pitfall}

Hajdu and Tijdeman, \emph{Algebraic aspects of discrete tomography},
Theorem 1 and its proof, describe planar switching components by
divisibility by products of direction binomials. Their argument uses
prime factors in a unique factorization domain to characterize zero
projections. It does not claim that arbitrary projection values are
compatible.
\url{https://math.unideb.hu/sites/default/files/inline-files/disctom8.pdf}

Tijdeman, \emph{Two applications of number theory to discrete tomography},
Theorem 3.1 as displayed, states a polynomial Chinese-remainder
assertion using pairwise relatively prime bivariate polynomials.
Gcd-one is insufficient for arbitrary residues. For instance,
$P_1=x$, $P_2=y$, $Q_1=0$, and $Q_2=1$ have no simultaneous
solution in $\mathbb R[x,y]$: evaluating the two congruences at
$(0,0)$ would give both $R(0,0)=0$ and $R(0,0)=1$.
Pairwise comaximal ideals, a stronger condition, give the usual CRT.
No arbitrary-residue conclusion from the displayed version is used
in this note. This counterexample does not contradict the earlier
zero-projection divisibility theorem.
\url{https://hrj.episciences.org/5111/pdf}

Davis and Sagan, \emph{Pattern-avoiding polytopes}, Proposition 3.4,
give the rectangular parallelotope of permutations avoiding
$132$ and $312$. Its orthogonal edge directions are the same
nested directions used by the triangular Boolean permutation
construction, up to coordinate conventions. The polytope result
does not provide arbitrary signs for a reciprocal readout.
\url{https://arxiv.org/pdf/1609.01782}

\section{Subsequent obstruction}

The kernel identities and finite compilers above remain valid.
The later proof in \texttt{laplace-sign-obstruction.tex} rules out
universality of the uniform-sign Mellin construction at
$H=O(2^n/n^2)$, even with arbitrarily large exponents. Evaluating
at $t=1/2$ gives an exponential score on the Boolean cube. It has
a polynomial parameterization of degree at most $n+1$, so its
strict-sign count is too small. Thus the sparse uniform-sign
representation is no longer treated as an open construction lemma.

\end{document}
